%% Generated by Sphinx.
\def\sphinxdocclass{report}
\documentclass[letterpaper,10pt,polish]{sphinxmanual}
\ifdefined\pdfpxdimen
   \let\sphinxpxdimen\pdfpxdimen\else\newdimen\sphinxpxdimen
\fi \sphinxpxdimen=.75bp\relax
\ifdefined\pdfimageresolution
    \pdfimageresolution= \numexpr \dimexpr1in\relax/\sphinxpxdimen\relax
\fi
%% let collapsible pdf bookmarks panel have high depth per default
\PassOptionsToPackage{bookmarksdepth=5}{hyperref}

\PassOptionsToPackage{booktabs}{sphinx}
\PassOptionsToPackage{colorrows}{sphinx}

\PassOptionsToPackage{warn}{textcomp}
\usepackage[utf8]{inputenc}
\ifdefined\DeclareUnicodeCharacter
% support both utf8 and utf8x syntaxes
  \ifdefined\DeclareUnicodeCharacterAsOptional
    \def\sphinxDUC#1{\DeclareUnicodeCharacter{"#1}}
  \else
    \let\sphinxDUC\DeclareUnicodeCharacter
  \fi
  \sphinxDUC{00A0}{\nobreakspace}
  \sphinxDUC{2500}{\sphinxunichar{2500}}
  \sphinxDUC{2502}{\sphinxunichar{2502}}
  \sphinxDUC{2514}{\sphinxunichar{2514}}
  \sphinxDUC{251C}{\sphinxunichar{251C}}
  \sphinxDUC{2572}{\textbackslash}
\fi
\usepackage{cmap}
\usepackage[T1]{fontenc}
\usepackage{amsmath,amssymb,amstext}
\usepackage{babel}



\usepackage{tgtermes}
\usepackage{tgheros}
\renewcommand{\ttdefault}{txtt}



\usepackage[Sonny]{fncychap}
\ChNameVar{\Large\normalfont\sffamily}
\ChTitleVar{\Large\normalfont\sffamily}
\usepackage{sphinx}

\fvset{fontsize=auto}
\usepackage{geometry}


% Include hyperref last.
\usepackage{hyperref}
% Fix anchor placement for figures with captions.
\usepackage{hypcap}% it must be loaded after hyperref.
% Set up styles of URL: it should be placed after hyperref.
\urlstyle{same}

\addto\captionspolish{\renewcommand{\contentsname}{Spis treści:}}

\usepackage{sphinxmessages}
\setcounter{tocdepth}{1}



\title{Sprawozdanie z laboratorium}
\date{15 cze 2026}
\release{0.0.1}
\author{Michał Bednarczyk}
\newcommand{\sphinxlogo}{\vbox{}}
\renewcommand{\releasename}{Wydanie}
\makeindex
\begin{document}

\ifdefined\shorthandoff
  \ifnum\catcode`\=\string=\active\shorthandoff{=}\fi
  \ifnum\catcode`\"=\active\shorthandoff{"}\fi
\fi

\pagestyle{empty}
\sphinxmaketitle
\pagestyle{plain}
\sphinxtableofcontents
\pagestyle{normal}
\phantomsection\label{\detokenize{index::doc}}


\sphinxAtStartPar
Add your content using \sphinxcode{\sphinxupquote{reStructuredText}} syntax. See the
\sphinxhref{https://www.sphinx-doc.org/en/master/usage/restructuredtext/index.html}{reStructuredText}
documentation for details.

\sphinxstepscope


\chapter{Wprowadzenie}
\label{\detokenize{rozdzial_1/index:wprowadzenie}}\label{\detokenize{rozdzial_1/index::doc}}
\sphinxstepscope


\chapter{Badania literaturowe}
\label{\detokenize{rozdzial_2/index:badania-literaturowe}}\label{\detokenize{rozdzial_2/index::doc}}
\sphinxstepscope


\section{Kopie zapasowe i odzyskiwanie danych}
\label{\detokenize{rozdzial_2/rozdzial_28/index:kopie-zapasowe-i-odzyskiwanie-danych}}\label{\detokenize{rozdzial_2/rozdzial_28/index::doc}}

\subsection{Wstęp}
\label{\detokenize{rozdzial_2/rozdzial_28/index:wstep}}
\sphinxAtStartPar
Kopie zapasowe w PostgreSQL można podzielić na dwa główne nurty: kopie logiczne wykonywane narzędziami \sphinxcode{\sphinxupquote{pg\_dump}} i \sphinxcode{\sphinxupquote{pg\_dumpall}} {[}cite: 3{]} oraz kopie fizyczne całego klastra wykonywane m.in. przez \sphinxcode{\sphinxupquote{pg\_basebackup}} wraz z archiwizacją WAL dla odzyskiwania punktowego PITR (Point\sphinxhyphen{}in\sphinxhyphen{}Time Recovery){[}cite: 3, 4{]}. Kopie logiczne są wygodne do odtwarzania pojedynczych obiektów i pojedynczych baz, natomiast kopie fizyczne są podstawą pełnego odtworzenia instancji, tablespaces i odzyskiwania po awariach{[}cite: 5{]}.


\subsection{Tworzenie kopii zapasowej całego systemu PostgreSQL \sphinxhyphen{} mechanizmy wbudowane}
\label{\detokenize{rozdzial_2/rozdzial_28/index:tworzenie-kopii-zapasowej-calego-systemu-postgresql-mechanizmy-wbudowane}}
\sphinxAtStartPar
Najprostszym wbudowanym sposobem wykonania kopii całego klastra jest \sphinxcode{\sphinxupquote{pg\_dumpall}}, które eksportuje wszystkie bazy w klastrze oraz obiekty globalne wspólne dla całej instancji, takie jak role i przestrzenie tabel{[}cite: 7{]}. Taki zrzut ma postać tekstowego skryptu SQL i odtwarza się go zwykle przez \sphinxcode{\sphinxupquote{psql}}, ale metoda ta jest logiczna, więc nie daje możliwości odzyskiwania punktowego przez WAL{[}cite: 8{]}.

\sphinxAtStartPar
Drugim, ważniejszym z punktu widzenia odporności na awarie mechanizmem jest \sphinxcode{\sphinxupquote{pg\_basebackup}}, które tworzy fizyczną kopię działającego klastra bez blokowania normalnej pracy klientów{[}cite: 10{]}. Narzędzie to wykonuje kopię całego klastra, a nie pojedynczych baz czy tabel, może pracować w formacie katalogowym lub \sphinxcode{\sphinxupquote{tar}} i może dołączać wymagane pliki WAL metodą stream, fetch albo none{[}cite: 11{]}. Dokumentacja podkreśla też, że taka kopia może służyć zarówno do PITR, jak i do zbudowania serwera standby{[}cite: 12{]}.

\sphinxAtStartPar
Aby pełna kopia fizyczna była użyteczna w scenariuszu odtwarzania do wybranego momentu, należy włączyć archiwizację WAL przez ustawienie co najmniej \sphinxcode{\sphinxupquote{wal\_level = replica}}, \sphinxcode{\sphinxupquote{archive\_mode = on}} oraz poprawnego \sphinxcode{\sphinxupquote{archive\_command}} albo {\color{red}\bfseries{}\textasciigrave{}\textasciigrave{}}archive\_library\textasciigrave{}\textasciigrave{}{[}cite: 13{]}. PostgreSQL zapisuje każdą zmianę w logu WAL, a po odtworzeniu kopii bazowej można odtworzyć kolejne segmenty WAL, aby dojść do stanu bieżącego lub do wskazanego punktu w czasie{[}cite: 14{]}.

\sphinxAtStartPar
Przykładowe polecenia:

\begin{sphinxVerbatim}[commandchars=\\\{\}]
\PYG{c+c1}{\PYGZsh{} logiczna kopia całego klastra}
pg\PYGZus{}dumpall\PYG{+w}{ }\PYGZhy{}U\PYG{+w}{ }postgres\PYG{+w}{ }\PYGZhy{}\PYGZhy{}clean\PYG{+w}{ }\PYGZhy{}\PYGZhy{}file\PYG{o}{=}cluster.sql

\PYG{c+c1}{\PYGZsh{} fizyczna kopia całego klastra z dołączeniem WAL}
pg\PYGZus{}basebackup\PYG{+w}{ }\PYGZhy{}h\PYG{+w}{ }dbserver\PYG{+w}{ }\PYGZhy{}D\PYG{+w}{ }/backup/base\PYG{+w}{ }\PYGZhy{}Fp\PYG{+w}{ }\PYGZhy{}X\PYG{+w}{ }stream\PYG{+w}{ }\PYGZhy{}P
\end{sphinxVerbatim}

\sphinxAtStartPar
W praktyce do dokumentacji laboratoryjnej warto zaznaczyć, że \sphinxcode{\sphinxupquote{pg\_dumpall}} lepiej nadaje się do migracji i prostych kopii administracyjnych, a \sphinxcode{\sphinxupquote{pg\_basebackup}} z archiwizacją WAL do scenariuszy disaster recovery{[}cite: 21{]}.


\subsection{Tworzenie kopii zapasowych poszczególnych baz danych \sphinxhyphen{} mechanizmy wbudowane}
\label{\detokenize{rozdzial_2/rozdzial_28/index:tworzenie-kopii-zapasowych-poszczegolnych-baz-danych-mechanizmy-wbudowane}}
\sphinxAtStartPar
Do wykonania kopii pojedynczej bazy służy \sphinxcode{\sphinxupquote{pg\_dump}}, które tworzy spójny zrzut nawet wtedy, gdy baza jest równocześnie używana przez innych użytkowników, i nie blokuje zwykłych operacji odczytu ani zapisu{[}cite: 24{]}. Narzędzie obsługuje format tekstowy, custom, directory oraz tar, przy czym formaty archiwalne współpracują z \sphinxcode{\sphinxupquote{pg\_restore}} i pozwalają na selektywne odtwarzanie obiektów{[}cite: 24{]}.

\sphinxAtStartPar
Istotne jest to, że \sphinxcode{\sphinxupquote{pg\_dump}} wykonuje kopię tylko jednej bazy danych {[}cite: 26{]}; do eksportu całego klastra lub obiektów globalnych należy użyć \sphinxcode{\sphinxupquote{pg\_dumpall\textasciigrave{}\textasciigrave{}{[}cite: 27{]}. Format custom (}}\sphinxhyphen{}Fc\textasciigrave{}\textasciigrave{}) i directory (\sphinxcode{\sphinxupquote{\sphinxhyphen{}Fd}}) są najbardziej elastyczne, bo pozwalają wybierać odtwarzane elementy, zmieniać kolejność odtwarzania, a w części scenariuszy także korzystać z odtwarzania równoległego{[}cite: 27{]}.

\sphinxAtStartPar
Przykładowe polecenia:

\begin{sphinxVerbatim}[commandchars=\\\{\}]
\PYG{c+c1}{\PYGZsh{} kopia pojedynczej bazy w formacie custom}
pg\PYGZus{}dump\PYG{+w}{ }\PYGZhy{}U\PYG{+w}{ }postgres\PYG{+w}{ }\PYGZhy{}d\PYG{+w}{ }labdb\PYG{+w}{ }\PYGZhy{}Fc\PYG{+w}{ }\PYGZhy{}f\PYG{+w}{ }/backup/labdb.dump

\PYG{c+c1}{\PYGZsh{} kopia pojedynczej bazy jako skrypt SQL}
pg\PYGZus{}dump\PYG{+w}{ }\PYGZhy{}U\PYG{+w}{ }postgres\PYG{+w}{ }\PYGZhy{}d\PYG{+w}{ }labdb\PYG{+w}{ }\PYGZhy{}Fp\PYG{+w}{ }\PYGZhy{}f\PYG{+w}{ }/backup/labdb.sql

\PYG{c+c1}{\PYGZsh{} odtworzenie archiwum custom}
pg\PYGZus{}restore\PYG{+w}{ }\PYGZhy{}d\PYG{+w}{ }labdb\PYGZus{}restored\PYG{+w}{ }/backup/labdb.dump
\end{sphinxVerbatim}

\sphinxAtStartPar
Jeżeli celem jest ochrona wybranych schematów lub tabel, \sphinxcode{\sphinxupquote{pg\_dump}} pozwala zawęzić zakres eksportu przez opcje takie jak \sphinxcode{\sphinxupquote{\sphinxhyphen{}n}} dla schematu i \sphinxcode{\sphinxupquote{\sphinxhyphen{}t}} dla tabeli{[}cite: 36{]}. Trzeba jednak zaznaczyć, że wybiórczy zrzut nie gwarantuje kompletności zależności wszystkich obiektów w nowym, pustym środowisku, jeśli pominięto elementy, od których zależy odtwarzany obiekt{[}cite: 37{]}.


\subsection{Odzyskiwanie usuniętych lub uszkodzonych danych}
\label{\detokenize{rozdzial_2/rozdzial_28/index:odzyskiwanie-usunietych-lub-uszkodzonych-danych}}
\sphinxAtStartPar
Sposób odzyskiwania zależy od tego, czy problem dotyczy pojedynczego obiektu logicznego, całej bazy, przestrzeni tabel, czy fizycznego uszkodzenia danych{[}cite: 40{]}. PostgreSQL rozróżnia w praktyce odzyskiwanie logiczne z plików dump oraz odzyskiwanie fizyczne z kopii bazowej i archiwum WAL{[}cite: 41{]}.


\subsubsection{Odtwarzanie tabel i innych obiektów logicznych}
\label{\detokenize{rozdzial_2/rozdzial_28/index:odtwarzanie-tabel-i-innych-obiektow-logicznych}}
\sphinxAtStartPar
Jeżeli wcześniej wykonano kopię \sphinxcode{\sphinxupquote{pg\_dump}} w formacie archiwalnym, można odtworzyć pojedynczą tabelę lub inny obiekt selektywnie przy pomocy \sphinxcode{\sphinxupquote{pg\_restore}}, bez przywracania całej bazy{[}cite: 43{]}. To jest najprostszy wariant odzyskania przypadkowo usuniętej tabeli, o ile odpowiedni obiekt znajdował się w logicznej kopii zapasowej{[}cite: 43{]}.

\sphinxAtStartPar
Przykład:

\begin{sphinxVerbatim}[commandchars=\\\{\}]
pg\PYGZus{}restore\PYG{+w}{ }\PYGZhy{}d\PYG{+w}{ }labdb\PYG{+w}{ }\PYGZhy{}t\PYG{+w}{ }public.wyniki\PYG{+w}{ }/backup/labdb.dump
\end{sphinxVerbatim}


\subsubsection{Odtwarzanie usuniętej bazy danych}
\label{\detokenize{rozdzial_2/rozdzial_28/index:odtwarzanie-usunietej-bazy-danych}}
\sphinxAtStartPar
Usuniętą bazę można najłatwiej odtworzyć z wcześniejszego zrzutu \sphinxcode{\sphinxupquote{pg\_dump}} tej konkretnej bazy albo z \sphinxcode{\sphinxupquote{pg\_dumpall}}, jeżeli wykonywano kopię całego klastra{[}cite: 48{]}. W przypadku wymogu odtworzenia dokładnie do chwili sprzed usunięcia, potrzebna jest kopia fizyczna oraz archiwum WAL, ponieważ same zrzuty logiczne nie są częścią rozwiązania continuous archiving i nie wspierają PITR{[}cite: 49{]}.


\subsubsection{Odtwarzanie przestrzeni tabel i pełnego klastra}
\label{\detokenize{rozdzial_2/rozdzial_28/index:odtwarzanie-przestrzeni-tabel-i-pelnego-klastra}}
\sphinxAtStartPar
Przestrzenie tabel są obiektami globalnymi klastra, dlatego ich definicje są uwzględniane przez \sphinxcode{\sphinxupquote{pg\_dumpall}}, a przy kopiach fizycznych \sphinxcode{\sphinxupquote{pg\_basebackup}} zachowuje także ich strukturę i mapowanie{[}cite: 52{]}. Dokumentacja \sphinxcode{\sphinxupquote{pg\_basebackup}} wskazuje, że przy pracy z tablespaces można użyć mapowania \sphinxcode{\sphinxupquote{\sphinxhyphen{}\sphinxhyphen{}tablespace\sphinxhyphen{}mapping}}, a przy odtwarzaniu trzeba zweryfikować poprawność dowiązań symbolicznych i lokalizacji danych{[}cite: 52{]}.

\sphinxAtStartPar
Przy fizycznym uszkodzeniu danych, katalogu PGDATA albo tabelpaces standardowa procedura obejmuje odtworzenie kopii bazowej, usunięcie starych plików danych, odtworzenie katalogów danych i tablespaces, skonfigurowanie \sphinxcode{\sphinxupquote{restore\_command}}, utworzenie pliku \sphinxcode{\sphinxupquote{recovery.signal}}, a następnie ponowne uruchomienie serwera w trybie recovery{[}cite: 54{]}. PostgreSQL podczas odtwarzania odczytuje wymagane segmenty WAL i może dojść do stanu bieżącego albo zatrzymać się na wybranym punkcie odzyskiwania{[}cite: 55{]}.


\subsubsection{PITR i odzyskiwanie stanu sprzed błędu}
\label{\detokenize{rozdzial_2/rozdzial_28/index:pitr-i-odzyskiwanie-stanu-sprzed-bledu}}
\sphinxAtStartPar
Najważniejszym mechanizmem odzyskiwania po usunięciu danych operacyjnych jest PITR, czyli przywrócenie systemu do chwili sprzed błędnej operacji, na przykład \sphinxcode{\sphinxupquote{DROP TABLE}} lub \sphinxcode{\sphinxupquote{DROP DATABASE\textasciigrave{}\textasciigrave{}{[}cite: 58{]}. W takim scenariuszu odtwarza się kopię bazową, konfiguruje \textasciigrave{}\textasciigrave{}restore\_command}}, a następnie ustawia cel odzyskiwania, np. przez czas, nazwany punkt przywracania lub identyfikator transakcji{[}cite: 59, 60{]}.

\sphinxAtStartPar
To rozwiązanie ma jednak ważne ograniczenie: nie służy do odzyskania pojedynczej tabeli „w miejscu” bez wpływu na resztę systemu, lecz do odtworzenia całego klastra do wcześniejszego stanu{[}cite: 61{]}. Dlatego w praktyce laboratoryjnej często łączy się oba podejścia: regularne \sphinxcode{\sphinxupquote{pg\_dump}} dla obiektów logicznych i \sphinxcode{\sphinxupquote{pg\_basebackup}} z archiwizacją WAL dla pełnego recovery{[}cite: 61{]}.


\subsection{Dedykowane oprogramowanie i skrypty zewnętrzne do automatyzacji}
\label{\detokenize{rozdzial_2/rozdzial_28/index:dedykowane-oprogramowanie-i-skrypty-zewnetrzne-do-automatyzacji}}
\sphinxAtStartPar
Wbudowane narzędzia PostgreSQL są wystarczające do ręcznego wykonywania kopii, ale w środowisku produkcyjnym często wykorzystuje się oprogramowanie automatyzujące harmonogramy, retencję, walidację i odtwarzanie{[}cite: 63{]}. Dwa najczęściej wskazywane rozwiązania to Barman i pgBackRest, przy czym dostępna dokumentacja Barman bezpośrednio opisuje obsługę pełnych i przyrostowych kopii, archiwizacji WAL oraz automatycznego uruchamiania backupów{[}cite: 64{]}.

\sphinxAtStartPar
Barman wykonuje kopie całego serwera PostgreSQL i wymaga poprawnej archiwizacji WAL, oferując różne metody kopii, w tym streaming przez \sphinxcode{\sphinxupquote{pg\_basebackup}}, \sphinxcode{\sphinxupquote{rsync}} oraz backupy snapshotowe w chmurze{[}cite: 66, 67{]}. Dokumentacja podaje też wsparcie dla backupów przyrostowych, limitowania pasma, kompresji i pracy z harmonogramem przez zewnętrzne mechanizmy systemowe, np. cron{[}cite: 67, 68{]}.

\sphinxAtStartPar
Przykładowe zastosowania narzędzi zewnętrznych:
\begin{itemize}
\item {} 
\sphinxAtStartPar
\sphinxstylestrong{Barman:} centralne repozytorium backupów, archiwizacja WAL, backupy pełne i przyrostowe, odzyskiwanie całych instancji{[}cite: 70{]}.

\item {} 
\sphinxAtStartPar
\sphinxstylestrong{pgBackRest:} popularne narzędzie do wydajnych backupów i retencji, często używane w środowiskach HA oraz większych instalacjach PostgreSQL{[}cite: 71{]}.

\item {} 
\sphinxAtStartPar
\sphinxstylestrong{Własne skrypty Bash/PowerShell:} wywołanie \sphinxcode{\sphinxupquote{pg\_dump}}, \sphinxcode{\sphinxupquote{pg\_dumpall}} lub \sphinxcode{\sphinxupquote{pg\_basebackup}}, kompresja plików, rotacja katalogów, logowanie i wysyłka wyników przez e\sphinxhyphen{}mail lub do systemu monitoringu{[}cite: 72{]}.

\end{itemize}

\sphinxAtStartPar
Przykładowy prosty skrypt automatyzujący kopię jednej bazy:

\begin{sphinxVerbatim}[commandchars=\\\{\}]
\PYG{c+ch}{\PYGZsh{}!/bin/bash}
\PYG{n+nv}{DATA}\PYG{o}{=}\PYG{k}{\PYGZdl{}(}date\PYG{+w}{ }+\PYGZpc{}F\PYGZus{}\PYGZpc{}H\PYGZhy{}\PYGZpc{}M\PYG{k}{)}
\PYG{n+nv}{KATALOG}\PYG{o}{=}/backup/postgres
\PYG{n+nv}{BAZA}\PYG{o}{=}labdb

mkdir\PYG{+w}{ }\PYGZhy{}p\PYG{+w}{ }\PYG{l+s+s2}{\PYGZdq{}}\PYG{n+nv}{\PYGZdl{}KATALOG}\PYG{l+s+s2}{\PYGZdq{}}
pg\PYGZus{}dump\PYG{+w}{ }\PYGZhy{}U\PYG{+w}{ }postgres\PYG{+w}{ }\PYGZhy{}d\PYG{+w}{ }\PYG{l+s+s2}{\PYGZdq{}}\PYG{n+nv}{\PYGZdl{}BAZA}\PYG{l+s+s2}{\PYGZdq{}}\PYG{+w}{ }\PYGZhy{}Fc\PYG{+w}{ }\PYGZhy{}f\PYG{+w}{ }\PYG{l+s+s2}{\PYGZdq{}}\PYG{n+nv}{\PYGZdl{}KATALOG}\PYG{l+s+s2}{/}\PYG{l+s+si}{\PYGZdl{}\PYGZob{}}\PYG{n+nv}{BAZA}\PYG{l+s+si}{\PYGZcb{}}\PYG{l+s+s2}{\PYGZus{}}\PYG{l+s+si}{\PYGZdl{}\PYGZob{}}\PYG{n+nv}{DATA}\PYG{l+s+si}{\PYGZcb{}}\PYG{l+s+s2}{.dump}\PYG{l+s+s2}{\PYGZdq{}}
find\PYG{+w}{ }\PYG{l+s+s2}{\PYGZdq{}}\PYG{n+nv}{\PYGZdl{}KATALOG}\PYG{l+s+s2}{\PYGZdq{}}\PYG{+w}{ }\PYGZhy{}type\PYG{+w}{ }f\PYG{+w}{ }\PYGZhy{}name\PYG{+w}{ }\PYG{l+s+s2}{\PYGZdq{}}\PYG{l+s+si}{\PYGZdl{}\PYGZob{}}\PYG{n+nv}{BAZA}\PYG{l+s+si}{\PYGZcb{}}\PYG{l+s+s2}{\PYGZus{}*.dump}\PYG{l+s+s2}{\PYGZdq{}}\PYG{+w}{ }\PYGZhy{}mtime\PYG{+w}{ }+7\PYG{+w}{ }\PYGZhy{}delete
\end{sphinxVerbatim}

\sphinxAtStartPar
Taki skrypt jest prosty, ale nie zapewnia pełnego disaster recovery, bo nie obejmuje archiwizacji WAL i nie chroni całego klastra{[}cite: 81{]}. Z tego powodu dla systemów krytycznych bardziej właściwe są narzędzia klasy Barman lub pgBackRest, które porządkują pełny proces backupu, retencji, testów i odtwarzania{[}cite: 81, 82{]}.


\subsection{Wnioski}
\label{\detokenize{rozdzial_2/rozdzial_28/index:wnioski}}
\sphinxAtStartPar
W rozdziale warto wyraźnie rozróżnić kopie logiczne i fizyczne, bo odpowiadają one na inne potrzeby eksploatacyjne{[}cite: 85{]}. \sphinxcode{\sphinxupquote{pg\_dump}} i \sphinxcode{\sphinxupquote{pg\_dumpall}} są najlepsze do prostych eksportów i selektywnego odtwarzania, natomiast \sphinxcode{\sphinxupquote{pg\_basebackup}} z archiwizacją WAL jest podstawą odzyskiwania po poważnej awarii i realizacji PITR{[}cite: 86{]}.

\sphinxAtStartPar
Dobrą praktyką jest także zapisanie, że skuteczna strategia backupu nie kończy się na wykonaniu kopii, ale obejmuje testy odtwarzania, kontrolę retencji, ochronę plików backupu oraz automatyzację procesu{[}cite: 88{]}. W laboratorium można więc pokazać zarówno proste polecenia wbudowane, jak i przykład narzędzia zewnętrznego, które organizuje cały cykl tworzenia i odtwarzania kopii zapasowych{[}cite: 88, 89{]}.

\sphinxstepscope


\chapter{Dokumentacja bazy danych \textendash{} Sklep elektroniczny}
\label{\detokenize{rozdzial_3/index:dokumentacja-bazy-danych-sklep-elektroniczny}}\label{\detokenize{rozdzial_3/index::doc}}
\sphinxAtStartPar
\sphinxstylestrong{Autorzy:} Norbert Antonovitch, Michał Bednarczyk, Jan Balazs de Borbatviz


\section{Wybór zagadnienia i identyfikacja danych}
\label{\detokenize{rozdzial_3/index:wybor-zagadnienia-i-identyfikacja-danych}}
\sphinxAtStartPar
Tematem projektu jest podstawowa baza danych obsługująca sklep internetowy ze sprzętem elektronicznym.
Baza danych zawiera następujące grupy danych:
\begin{itemize}
\item {} 
\sphinxAtStartPar
\sphinxstylestrong{Klienci:} unikalny identyfikator, imię, nazwisko, adres e\sphinxhyphen{}mail, numer telefonu oraz podstawowe dane adresowe do wysyłki: ulica, miasto, kod pocztowy.

\item {} 
\sphinxAtStartPar
\sphinxstylestrong{Produkty:} unikalny identyfikator, nazwa sprzętu, cena jednostkowa oraz przypisana kategoria.

\item {} 
\sphinxAtStartPar
\sphinxstylestrong{Kategorie:} identyfikator kategorii oraz jej nazwa, na przykład Laptopy lub Smartfony.

\item {} 
\sphinxAtStartPar
\sphinxstylestrong{Zamówienia:} identyfikator zamówienia, data złożenia, status, na przykład nowe lub zrealizowane, oraz powiązanie z konkretnym klientem.

\item {} 
\sphinxAtStartPar
\sphinxstylestrong{Szczegóły zamówienia:} pozycje przypisane do konkretnego zamówienia, określające jaki produkt został kupiony, w jakiej ilości oraz po jakiej cenie.

\end{itemize}


\section{Opis procesów i towarzyszących im danych}
\label{\detokenize{rozdzial_3/index:opis-procesow-i-towarzyszacych-im-danych}}\begin{enumerate}
\sphinxsetlistlabels{\arabic}{enumi}{enumii}{}{.}%
\item {} 
\sphinxAtStartPar
\sphinxstylestrong{Zarządzanie produktami:} sklep oferuje asortyment przypisany do konkretnych kategorii, na przykład laptopy i smartfony. Każdy produkt ma przypisaną nazwę, cenę oraz identyfikator producenta.

\item {} 
\sphinxAtStartPar
\sphinxstylestrong{Zarządzanie użytkownikami:} klienci rejestrują się w systemie, podając podstawowe dane, takie jak imię, nazwisko i e\sphinxhyphen{}mail, oraz pojedynczy główny adres zamieszkania wykorzystywany do wysyłki.

\item {} 
\sphinxAtStartPar
\sphinxstylestrong{Realizacja zamówień:} zalogowany klient składa zamówienie na wybrane produkty. System tworzy nagłówek zamówienia zawierający datę i status oraz przypisuje do niego poszczególne produkty wraz z ich ilością.

\end{enumerate}


\section{Prototypowy plik JSON}
\label{\detokenize{rozdzial_3/index:prototypowy-plik-json}}
\sphinxAtStartPar
Poniżej przedstawiono prototypowy plik w formacie JSON, zawierający jeden pełny dokument reprezentujący złożone zamówienie. Pozwala to zweryfikować kompletność gromadzonych i przetwarzanych informacji przed przejściem do modelowania konceptualnego.

\begin{sphinxVerbatim}[commandchars=\\\{\}]
\PYG{p}{\PYGZob{}}
\PYG{+w}{  }\PYG{n+nt}{\PYGZdq{}id\PYGZus{}zamowienia\PYGZdq{}}\PYG{p}{:}\PYG{+w}{ }\PYG{l+m+mi}{1}\PYG{p}{,}
\PYG{+w}{  }\PYG{n+nt}{\PYGZdq{}data\PYGZus{}zlozenia\PYGZdq{}}\PYG{p}{:}\PYG{+w}{ }\PYG{l+s+s2}{\PYGZdq{}2026\PYGZhy{}05\PYGZhy{}14T10:15:30Z\PYGZdq{}}\PYG{p}{,}
\PYG{+w}{  }\PYG{n+nt}{\PYGZdq{}status\PYGZdq{}}\PYG{p}{:}\PYG{+w}{ }\PYG{l+s+s2}{\PYGZdq{}Nowe\PYGZdq{}}\PYG{p}{,}
\PYG{+w}{  }\PYG{n+nt}{\PYGZdq{}klient\PYGZdq{}}\PYG{p}{:}\PYG{+w}{ }\PYG{p}{\PYGZob{}}
\PYG{+w}{    }\PYG{n+nt}{\PYGZdq{}id\PYGZus{}klienta\PYGZdq{}}\PYG{p}{:}\PYG{+w}{ }\PYG{l+m+mi}{4096}\PYG{p}{,}
\PYG{+w}{    }\PYG{n+nt}{\PYGZdq{}imie\PYGZdq{}}\PYG{p}{:}\PYG{+w}{ }\PYG{l+s+s2}{\PYGZdq{}Jan\PYGZdq{}}\PYG{p}{,}
\PYG{+w}{    }\PYG{n+nt}{\PYGZdq{}nazwisko\PYGZdq{}}\PYG{p}{:}\PYG{+w}{ }\PYG{l+s+s2}{\PYGZdq{}Kowalski\PYGZdq{}}\PYG{p}{,}
\PYG{+w}{    }\PYG{n+nt}{\PYGZdq{}email\PYGZdq{}}\PYG{p}{:}\PYG{+w}{ }\PYG{l+s+s2}{\PYGZdq{}student01@example.com\PYGZdq{}}\PYG{p}{,}
\PYG{+w}{    }\PYG{n+nt}{\PYGZdq{}telefon\PYGZdq{}}\PYG{p}{:}\PYG{+w}{ }\PYG{l+s+s2}{\PYGZdq{}123456789\PYGZdq{}}\PYG{p}{,}
\PYG{+w}{    }\PYG{n+nt}{\PYGZdq{}ulica\PYGZdq{}}\PYG{p}{:}\PYG{+w}{ }\PYG{l+s+s2}{\PYGZdq{}Armii Krajowej 78\PYGZdq{}}\PYG{p}{,}
\PYG{+w}{    }\PYG{n+nt}{\PYGZdq{}kod\PYGZus{}pocztowy\PYGZdq{}}\PYG{p}{:}\PYG{+w}{ }\PYG{l+s+s2}{\PYGZdq{}58\PYGZhy{}300\PYGZdq{}}\PYG{p}{,}
\PYG{+w}{    }\PYG{n+nt}{\PYGZdq{}miasto\PYGZdq{}}\PYG{p}{:}\PYG{+w}{ }\PYG{l+s+s2}{\PYGZdq{}Walbrzych\PYGZdq{}}
\PYG{+w}{  }\PYG{p}{\PYGZcb{},}
\PYG{+w}{  }\PYG{n+nt}{\PYGZdq{}pozycje\PYGZdq{}}\PYG{p}{:}\PYG{+w}{ }\PYG{p}{[}
\PYG{+w}{    }\PYG{p}{\PYGZob{}}
\PYG{+w}{      }\PYG{n+nt}{\PYGZdq{}id\PYGZus{}produktu\PYGZdq{}}\PYG{p}{:}\PYG{+w}{ }\PYG{l+m+mi}{10543}\PYG{p}{,}
\PYG{+w}{      }\PYG{n+nt}{\PYGZdq{}nazwa\PYGZdq{}}\PYG{p}{:}\PYG{+w}{ }\PYG{l+s+s2}{\PYGZdq{}Laptop ASUS\PYGZdq{}}\PYG{p}{,}
\PYG{+w}{      }\PYG{n+nt}{\PYGZdq{}kategoria\PYGZdq{}}\PYG{p}{:}\PYG{+w}{ }\PYG{l+s+s2}{\PYGZdq{}Laptopy\PYGZdq{}}\PYG{p}{,}
\PYG{+w}{      }\PYG{n+nt}{\PYGZdq{}ilosc\PYGZdq{}}\PYG{p}{:}\PYG{+w}{ }\PYG{l+m+mi}{1}\PYG{p}{,}
\PYG{+w}{      }\PYG{n+nt}{\PYGZdq{}cena\PYGZus{}historyczna\PYGZdq{}}\PYG{p}{:}\PYG{+w}{ }\PYG{l+m+mf}{6299.99}
\PYG{+w}{    }\PYG{p}{\PYGZcb{}}
\PYG{+w}{  }\PYG{p}{]}
\PYG{p}{\PYGZcb{}}
\end{sphinxVerbatim}


\subsection{Dobór typów danych}
\label{\detokenize{rozdzial_3/index:dobor-typow-danych}}

\begin{savenotes}\sphinxattablestart
\sphinxthistablewithglobalstyle
\centering
\sphinxcapstartof{table}
\sphinxthecaptionisattop
\sphinxcaption{Mapowanie typów danych dla SQLite i PostgreSQL}\label{\detokenize{rozdzial_3/index:id1}}
\sphinxaftertopcaption
\begin{tabulary}{\linewidth}[t]{TTT}
\sphinxtoprule
\sphinxstyletheadfamily 
\sphinxAtStartPar
Atrybut danych
&\sphinxstyletheadfamily 
\sphinxAtStartPar
SQLite
&\sphinxstyletheadfamily 
\sphinxAtStartPar
PostgreSQL
\\
\sphinxmidrule
\sphinxtableatstartofbodyhook
\sphinxAtStartPar
Identyfikatory (\sphinxcode{\sphinxupquote{id\_...}})
&
\sphinxAtStartPar
\sphinxcode{\sphinxupquote{INTEGER}}
&
\sphinxAtStartPar
\sphinxcode{\sphinxupquote{SERIAL}} lub \sphinxcode{\sphinxupquote{INTEGER}}
\\
\sphinxhline
\sphinxAtStartPar
Napisy krótkie (\sphinxcode{\sphinxupquote{imie}}, \sphinxcode{\sphinxupquote{nazwisko}}, \sphinxcode{\sphinxupquote{miasto}})
&
\sphinxAtStartPar
\sphinxcode{\sphinxupquote{TEXT}}
&
\sphinxAtStartPar
\sphinxcode{\sphinxupquote{VARCHAR(50)}}
\\
\sphinxhline
\sphinxAtStartPar
Napisy długie (\sphinxcode{\sphinxupquote{ulica}}, \sphinxcode{\sphinxupquote{email}}, \sphinxcode{\sphinxupquote{nazwa}})
&
\sphinxAtStartPar
\sphinxcode{\sphinxupquote{TEXT}}
&
\sphinxAtStartPar
\sphinxcode{\sphinxupquote{VARCHAR(255)}}
\\
\sphinxhline
\sphinxAtStartPar
Stałe ciągi znaków (\sphinxcode{\sphinxupquote{telefon}}, \sphinxcode{\sphinxupquote{kod\_pocztowy}})
&
\sphinxAtStartPar
\sphinxcode{\sphinxupquote{TEXT}}
&
\sphinxAtStartPar
\sphinxcode{\sphinxupquote{VARCHAR(15)}}
\\
\sphinxhline
\sphinxAtStartPar
Cena (\sphinxcode{\sphinxupquote{cena\_bazowa}}, \sphinxcode{\sphinxupquote{cena\_historyczna}})
&
\sphinxAtStartPar
\sphinxcode{\sphinxupquote{FLOAT}}
&
\sphinxAtStartPar
\sphinxcode{\sphinxupquote{NUMERIC(10,2)}}
\\
\sphinxhline
\sphinxAtStartPar
Ilość (\sphinxcode{\sphinxupquote{ilosc}})
&
\sphinxAtStartPar
\sphinxcode{\sphinxupquote{INTEGER}}
&
\sphinxAtStartPar
\sphinxcode{\sphinxupquote{SMALLINT}}
\\
\sphinxhline
\sphinxAtStartPar
Data i czas (\sphinxcode{\sphinxupquote{data\_zlozenia}})
&
\sphinxAtStartPar
\sphinxcode{\sphinxupquote{TEXT}}
&
\sphinxAtStartPar
\sphinxcode{\sphinxupquote{TIMESTAMP}}
\\
\sphinxbottomrule
\end{tabulary}
\sphinxtableafterendhook\par
\sphinxattableend\end{savenotes}


\section{Modelowanie konceptualne i logiczne}
\label{\detokenize{rozdzial_3/index:modelowanie-konceptualne-i-logiczne}}

\subsection{Model koncepcyjny}
\label{\detokenize{rozdzial_3/index:model-koncepcyjny}}
\sphinxAtStartPar
\sphinxstylestrong{1. Encje mocne}
\begin{itemize}
\item {} 
\sphinxAtStartPar
\sphinxstylestrong{Klient:} klucz główny (\sphinxcode{\sphinxupquote{id\_klienta}}), imię, nazwisko, e\sphinxhyphen{}mail, telefon, ulica, miasto, kod pocztowy.

\item {} 
\sphinxAtStartPar
\sphinxstylestrong{Produkt:} klucz główny (\sphinxcode{\sphinxupquote{id\_produktu}}), nazwa, \sphinxcode{\sphinxupquote{cena\_bazowa}}.

\item {} 
\sphinxAtStartPar
\sphinxstylestrong{Kategoria:} klucz główny (\sphinxcode{\sphinxupquote{id\_kategorii}}), nazwa.

\item {} 
\sphinxAtStartPar
\sphinxstylestrong{Zamówienie:} klucz główny (\sphinxcode{\sphinxupquote{id\_zamowienia}}), \sphinxcode{\sphinxupquote{data\_zlozenia}}, status.

\end{itemize}

\sphinxAtStartPar
\sphinxstylestrong{2. Encje słabe}
\begin{itemize}
\item {} 
\sphinxAtStartPar
\sphinxstylestrong{Szczegóły zamówienia:} klucz obcy (\sphinxcode{\sphinxupquote{id\_zamowienia}}), klucz obcy (\sphinxcode{\sphinxupquote{id\_produktu}}), ilość, \sphinxcode{\sphinxupquote{cena\_historyczna}}. Jej identyfikacja opiera się wyłącznie na kluczach obcych pochodzących z encji silnych.

\end{itemize}

\sphinxAtStartPar
\sphinxstylestrong{3. Opis związków}

\sphinxAtStartPar
Zdefiniowano następujące relacje biznesowe między encjami:
\begin{itemize}
\item {} 
\sphinxAtStartPar
\sphinxstylestrong{Klient \textendash{} Zamówienie (1:N):} klient składa zamówienie. Jeden klient może posiadać wiele zamówień, ale jedno konkretne zamówienie jest przypisane dokładnie do jednego klienta.

\item {} 
\sphinxAtStartPar
\sphinxstylestrong{Kategoria \textendash{} Produkt (1:N):} kategoria grupuje produkty. Kategoria może zawierać wiele produktów, ale dany produkt należy do jednej kategorii.

\item {} 
\sphinxAtStartPar
\sphinxstylestrong{Zamówienie \textendash{} Produkt (M:N):} zamówienie zawiera produkty. Pojedyncze zamówienie obejmuje wiele produktów, a dany produkt może pojawić się w wielu zamówieniach.

\end{itemize}

\sphinxAtStartPar
\sphinxstylestrong{4. Związki niepoprawne}

\sphinxAtStartPar
Wyeliminowano bezpośrednią relację między Klientem a Produktem, aby uniknąć pułapki połączeń. Bezpośrednie powiązanie gubi kontekst transakcji, na przykład datę i ilość zakupionego towaru, dlatego poprawna relacja musi zawsze przechodzić przez encję Zamówienie.


\subsection{Model koncepcyjny \textendash{} notacja Chena}
\label{\detokenize{rozdzial_3/index:model-koncepcyjny-notacja-chena}}\begin{quote}

\noindent\sphinxincludegraphics[width=0.800\linewidth]{{schemat_chen}.png}
\end{quote}


\subsection{Model logiczny i normalizacja}
\label{\detokenize{rozdzial_3/index:model-logiczny-i-normalizacja}}
\sphinxAtStartPar
Główne działania podjęte podczas normalizacji:
\begin{enumerate}
\sphinxsetlistlabels{\arabic}{enumi}{enumii}{}{.}%
\item {} 
\sphinxAtStartPar
\sphinxstylestrong{I Postać Normalna (1NF):} dane adresowe klienta, takie jak ulica, miasto i kod pocztowy, zostały wyodrębnione jako pojedyncze pola.

\item {} 
\sphinxAtStartPar
\sphinxstylestrong{II Postać Normalna (2NF):} zlikwidowano relację wiele\sphinxhyphen{}do\sphinxhyphen{}wielu (N:M) pomiędzy Zamówieniem a Produktem, wstawiając encję asocjacyjną \sphinxcode{\sphinxupquote{Szczegóły\_zamówienia}}. Jej klucz główny składa się z identyfikatora zamówienia oraz identyfikatora produktu.

\item {} 
\sphinxAtStartPar
\sphinxstylestrong{III Postać Normalna (3NF):} usunięto zależności przechodnie. Nazwa kategorii została wydzielona do osobnej tabeli Kategorie. W tabeli Produkty przechowywany jest jedynie klucz obcy do kategorii. Ponadto zapewniono zapisywanie historycznej ceny w \sphinxcode{\sphinxupquote{Szczegółach\_zamówienia}}, aby modyfikacja cennika nie fałszowała archiwalnych rachunków.

\end{enumerate}


\subsection{Schemat notacji IE po normalizacji}
\label{\detokenize{rozdzial_3/index:schemat-notacji-ie-po-normalizacji}}\begin{quote}

\noindent\sphinxincludegraphics[width=0.800\linewidth]{{schemat_ie}.png}
\end{quote}

\sphinxstepscope



\renewcommand{\indexname}{Indeks}
\printindex
\end{document}